% Part: first-order-logic
% Chapter: syntax
% Section: intro-syntax

\documentclass[../../../include/open-logic-section]{subfiles}

\begin{document}

\olfileid{fol}{syn}{itx}

\olsection{Inleiding}

Om de logica van de eerste orde te bestuderen, moeten we ook
onderzoeken hoe dit logische systeem werkt. Daarom leggen we eerst de
vormregels (syntaxis) en de betekenisregels (semantiek) van de
uitdrukkingen vast. De uitdrukkingen van de logica van de eerste orde
zijn termen en !!{formula}s. Termen worden opgebouwd uit
!!{variable}s, !!{constant}s en !!{function}s.
De !!^{formula}s worden weer uit andere onderdelen opgebouwd. Je
combineert !!{predicate}s met termen; dat levert de kleinste,
``atomaire'' !!{formula}s op. Daarna maak je met logische
verbindingstekens en kwantoren uit die atomaire !!{formula}s
ingewikkeldere formules.
Je kunt daarvoor verschillende vormregels kiezen. Dit boek kiest één
mogelijke set. Andere systemen kunnen andere symbolen gebruiken. Ze
kunnen ook andere verbindingstekens als basis nemen, en haakjes anders
of helemaal niet gebruiken. Bij de zogenoemde Poolse notatie zijn
bijvoorbeeld geen haakjes nodig.
Wat al die keuzes gemeen hebben, is dat de vormingsregels de
verzameling termen en !!{formula}s stap voor stap \emph{inductief}
vastleggen. Als de regels goed zijn opgesteld, kun je in wezen maar op
één manier nagaan hoe een uitdrukking is opgebouwd. Elke uitdrukking is
dan \emph{eenduidig leesbaar}. Daarom kunnen we ook haar betekenis stap
voor stap vastleggen, met hetzelfde soort inductieve definitie.
\end{document}
